% Polaris 稿件骨架（NeurIPS 2026，开发用简化样式）
% 分节占位注释（% POLARIS_SECTION: xxx ... % POLARIS_SECTION_END: xxx）之间的内容
% 由 AI 起草 / 协同编辑写入，请勿删除这些标记行。
\documentclass{article}
\usepackage{polaris_neurips2026}
\usepackage[utf8]{inputenc}
\usepackage{graphicx}
\usepackage{amsmath}
\usepackage{booktabs}
\usepackage{natbib}

\title{{{POLARIS_TITLE}}}
\author{Polaris Draft}

\begin{document}
\maketitle

\begin{abstract}
% POLARIS_SECTION: abstract
（待撰写 / to be drafted）
% POLARIS_SECTION_END: abstract
\end{abstract}

\section{Introduction}\label{sec:introduction}
% POLARIS_SECTION: introduction
（待撰写 / to be drafted）
% POLARIS_SECTION_END: introduction

\section{Related Work}\label{sec:related_work}
% POLARIS_SECTION: related_work
（待撰写 / to be drafted）
% POLARIS_SECTION_END: related_work

\section{Method}\label{sec:method}
% POLARIS_SECTION: method
（待撰写 / to be drafted）
% POLARIS_SECTION_END: method

\section{Experimental Setup}\label{sec:experimental_setup}
% POLARIS_SECTION: experimental_setup
（待撰写 / to be drafted）
% POLARIS_SECTION_END: experimental_setup

\section{Results}\label{sec:results}
% POLARIS_SECTION: results
（待撰写 / to be drafted）
% POLARIS_SECTION_END: results

\section{Conclusion}\label{sec:conclusion}
% POLARIS_SECTION: conclusion
（待撰写 / to be drafted）
% POLARIS_SECTION_END: conclusion

\bibliographystyle{plainnat}
\bibliography{references}

\end{document}
